% Pursuing Stacks (cotes 134-2 to 134-8): archive cote 134-2, site page 62 (PDF page 63).
% appendix (letters to Breen). Draft compared: ps1.tex:1454-1468.
% AI-transcribed by Claude Opus 5.5 (claude-opus-5-5, adaptive thinking, effort high), Message Batches API, batch ps1 (msgbatch_01Rv5hCaNTtUt4JHx1yqoHPL), 2026-11-12; not yet checked by eye.
\scanpage{134-2/62}
\mspage{-47 -}

\margin{(App.12)}(D) Let $X/k$ be a smooth scheme (for simplicity) \deleted{on}\inserted{over} a field $k$ , separated
and of finite type, of relative dimension $d$ , and $n$ an integer $>0$ .
If $n$ is prime to the characteristic and if $F$ is a sheaf of coefficients
on $X$ which is annihilated by $n$ , the global duality tells us that $Rf_!(F)$
and $Rf_*\,(R\,\mathit{Hom}(F,\mu_n^{\otimes d}))$ ( $\mu_n$ = sheaf of nth roots of unity = $\ker(G_m \xrightarrow{n} G_m)$)
are dual to each other with values in $(\mathbb{Z}/n\mathbb{Z})_k$ , for example $Rf_!(\mathbb{Z}/n\mathbb{Z})$
and $Rf_*(\mu_n^{\otimes d})$ , or $Rf_!(\mu_m^{\otimes})$ and $Rf_*(\mathbb{Z}/n\mathbb{Z})$ , are dual to each other - at
least with a shift of amplitude $2d$ in dimension. (As $\mathbb{Z}/n\mathbb{Z}$ is injective
over itself, this gives in fact perfect dualities
\[ R^if_!(F) \times R^{2d-i}f_*(\mathbb{R}\mathrm{hom}(F, \mu_n^{\otimes d})) \longrightarrow \mathbb{Z}/n\mathbb{Z}\deleted{)}. \hand{)} \]
If now one no longer assumes $n$ prime to the characteristic, for example $n$
is a power of $p$ = characteristic of $k > 0$), it seems that everything
collapses: to start with, one no longer knows (for $d > 1$) by what to replace
$\mu_n^{\otimes d}$ ...

The extraordinary miracle is that for $d = 1$ , i.e. $X$ a smooth curve,
everything continues to work perfectly, provided one states things with care!
The first verifications are made for example with $F = \mathbb{Z}/\deleted{\illegible{1}}\hand{p}\mathbb{Z}$\hand{,} $\mu_p$, \inserted{or} $\alpha_p$ , \deleted{\unsure{X|K}}\hand{with X}
complete - one finds it's O.K. by virtue essentially of the autoduality
of the jacobian. One can make these examples more sophisticated on taking
\textit{twisted} coefficients, and $X$ not complete - one convinces oneself this works
always! Simply, it is necessary to note that here the $R^if_*(F)$ , $R^if_!(F)$
have a ``continuous'' structure (they are essentially proalgebraic groups).
This corresponds to the well known phenomenon in class field theory that the
structure of $\pi_{1ab}$ of $X$ , when $X$ is not complete, is \textit{continuous} - hence
\margin{a/ \quad $p^n\mathbb{Z}$}\unsure{same} holds for $H^1(X, \mathbb{Z}/\deleted{\unsure{n}}\hand{p^n}\mathbb{Z})$ \hand{say}.

By the way, I point out for you that Serre once proposed (without ever
writing it down, I think) a theory of duality for \uline{commutative} \uline{unipotent}
algebraic groups, \textit{mod radical isogeny}, duality with values in $\mathbb{Q}/\mathbb{Z}$ (or
$\mathbb{Q}_p/\mathbb{Z}_p$) . He found that if (when $k$ is algebraically closed, say) $G$ is

\margin{62}
